% ISWC 2026 In-Use Track submission
% Target: https://iswc2026.semanticweb.org/#/calls/in-use
%
% ISWC 2026 In-Use Track CFP Key Dates:
%   Abstract submission due:    2 May 2026
%   Full paper submission due:  7 May 2026
%   Rebuttal:                   11 June - 18 June 2026
%   Notifications:              16 July 2026
%   Camera-ready papers due:    6 August 2026
%   Conference:                 27 - 29 October 2026
%   All deadlines 23:59 AoE
%
% Submission details:
%   Max 15 pages (excl. references and "Declaration of use of Generative AI")
%   Format: LNCS (Springer Lecture Notes in Computer Science)
%   Review: single-blind (not anonymous)
%   Submission via EasyChair (abstract pre-submission required)
%   Published: Springer LNCS proceedings
%
% Track Chairs:
%   Ernesto Jimenez-Ruiz, City St George's, University of London, UK
%   Oshani Seneviratne, Rensselaer Polytechnic Institute, USA
%   Contact: iswc2026-in-use@easychair.org
%
% Based on prior SWJ submission, now migrated to LNCS
% Original: extended version of Text2KG ESWC 2023 Workshop paper

\documentclass[runningheads]{llncs}

%% Packages
\usepackage{dcolumn}

%% Definitions
\newcolumntype{d}[1]{D{.}{.}{#1}}

%% Authors LaTex packages and macros
% enable utf-8 chars
\usepackage[utf8]{inputenc}

% according to error message
\usepackage{booktabs}

% render SVG files: https://tex.stackexchange.com/questions/122871/include-svg-images-with-the-svg-package
\usepackage{svg}

% for tables
\usepackage{tabularx}

% source code rendering
\usepackage{minted}

% hyperlinks
\usepackage{hyperref}
\usepackage{url}

% ORCID rendering provided by llncs.cls

% typographical quoting
\newcommand{\tquote}[1]{``#1''}
% footnotelink
\newcommand{\footnotelink}[2]{\href{#1}{#2}\footnote{#1}}
% define some macros/commands
\newcommand{\wdq}[2]{\footnotelink{https://www.wikidata.org/wiki/#2}{#1 (#2)}}
\newcommand{\wdp}[2]{\footnotelink{https://www.wikidata.org/wiki/Property:#2}{#1 (#2)}}

\newcommand{\wdpTable}[2]{\href{https://www.wikidata.org/wiki/Property:#2}{#1 (#2)}}
\newcommand{\wdqTable}[2]{\href{https://www.wikidata.org/wiki/#2}{#1 (#2)}}

%% Article Info


\title{Semantified CEUR-WS in Wikidata}

\author{Wolfgang Fahl\inst{1}\orcidID{0000-0002-0821-6995} \and
Tim Holzheim\inst{1,2}\orcidID{0000-0003-2533-6363} \and
Christoph Lange\inst{1,2}\orcidID{0000-0001-9879-3827} \and
Jerven Bolleman\inst{3}\orcidID{0000-0002-7449-1266} \and
Stefan Decker\inst{1,2}\orcidID{0000-0001-6324-7164}}

\institute{RWTH Aachen University, Computer Science i5, Aachen, Germany \and
Fraunhofer FIT, Sankt Augustin, Germany \and
SIB Swiss Institute of Bioinformatics, Switzerland}

\begin{document}
\authorrunning{Wolfgang Fahl et. al.}
\maketitle
\begin{abstract}
CEUR Workshop Proceedings (CEUR-WS) has been a non-profit publisher for scientific workshop and conference proceedings since 1995. For more than 30 years it has used HTML and PDF as the single source of truth for the metadata required by indexers such as DBLP and the German national library DNB. Scholars have long craved persistent identifiers for the core entities involved: event series, events, proceedings, papers, editors, authors, institutions, and publishers, with the goal of allowing semantic handling and querying with modern tools.

In this work we show how we enabled the semantification using Wikidata as the target knowledge graph infrastructure. The 2015 Semantic Publishing Challenge queries are finally permanently available with up-to-date metadata. The Scholia portal provides aspect-oriented and user-friendly access and is widely used.

We further point out the challenges of aligning the digital, social, and physical aspects of the ongoing transformation and outline a roadmap to use the available queries to bootstrap a new RDF-based metadata-first publishing approach that will mine metadata from proceedings matter and generate the HTML from it instead of the other way around.


\end{abstract}
\keywords{Knowledge Graph \and Linked Data \and Metadata Extraction \and Publishing \and Semantic Web \and Wikidata}

\section{Introduction}\label{sec:Introduction}
\begin{figure}
    \centering
    \includesvg[width=\linewidth]{images/CoreEntities.svg}
    \caption{Most relevant entities for scientific publishing (abstract schema). Source: CR research wiki~\cite{Fahl_CR_Wiki}.}
     \label{fig:core_entities}
\end{figure}

\begin{figure}
    \centering
    \includesvg[width=\linewidth]{images/SWAT4HCLS2024.svg}
    \caption{SWAT4HCLS2024 example graph walk}
    \label{fig:swat4hcls_2024_instance}
\end{figure}

The Traditional publishing workflows often handle the metadata of the needed core entities (as shown in Figure~\ref{fig:core_entities}) in an unFAIR way. Findable, Accessible, Interoperable, and Reusable (FAIR) ~\ref{sec:FAIR}, Single Point of Truth (SPoT), and Single Source of Truth (SSoT) principles are having low priority. Therefore the opportunities to collect the metadata during the different lifcecyle stages of an event ~\cite{FahlHOD23} are missed. As technical editorr of CEUR-WS we have been feeling this pain for over a decade. This work shows how the HTML/PDF SSoT of CEUR-WS at the ceur-ws.org domain (SPoT) has been used to create a Knowledge Graph representation of the relevant entities in Wikidata between 2022 and 2026 as a step in this direction thus allowing portals such as Scholia to present the RDF query results in an aspect-oriented and user-friendly way.

\subsection{CEUR Workshop Proceedings}
Since 1995 CEUR Workshop Proceedings (CEUR-WS) (https://ceur-ws.org/) has published more than 4,200 Volumes and over 90.000 papers as a free online service for publishing proceedings of scientific workshops (and smaller conferences) in computer science, operated by at team of unpaid volunteers.

Technically, all data and metadata of the proceedings are directly represented as a static website using HTML, PDF and HTTP and FTP protocols as provided by the workshop editors along the CEUR-WS “how to submit” guide~\cite{CEURHowtoSubmit}.

The metadata is available after some delay of weeks or month through indexing services such as dblp~\cite{Ley2009} and {K10plus}\cite{k10plus2019}\footnote{\url{https://dblp.org/}, \url{https://opac.k10plus.de}} and at the archive of the TIB \url{https://www.tib.eu/}.

After multiple unfinished attempts to make the metadata of the CEUR-WS platform available for semantic analysis and querying we report on the successful use of wikidata and scholia as the knowledge graph platform to provide FAIR metadata.

Doing so was harder than one would expect - the devil is in the details.

\subsection{Making CEUR-WS FAIR}\label{sec:FAIR}
% FAIR Data Introduction - why is this relevant? EOSC, NFDI, Open Access, Open Science
% FAIR Data as a basis for Knowledge Graphs
The FAIR principles~\cite{GOFAIR2019,Wilkinson2016} with Persistent Identifiers (PIDs) and rich metadata as core components are being adopted by an increasing number of publishing organizations.

To make CEUR-WS metadata fair a set of relevant queries should be supported, as derived from the original set of queries of the 2014 Semantic publishing challenge as outlined in Section~\ref{sec:SemantificationStory}.
The metadata should be provided following an ontology that is based on established ontologies.
The manual and automatic curation of entries should be possible with public access for all stakeholders, e.g., editors, authors, organizers, publishers and indexers.
The infrastructure should be stable and there should be sufficient trust in its long term availability
An open source non-profit infrastructure is preferred since this is also the mode of operation of CEUR-WS

Given the HTML/PDF/text input of CEUR-WS, the corresponding Knowledge Graph needs to be created and the single-point-of-truth computer readable metadata be separated from the different representations such as HTML so that the above requirements are fulfilled.

%The \footnotelink{https://cr.bitplan.com/index.php/Semantic\_Publishing\_Challenge}{original task} called for 20 such queries to be answerable.
%Formally we need to implement a function hat takes two Trees (IT,PT) as an Input and outputs a Graph G=(V,E).

%The IT Tree is the HTML Dom Tree of the CEUR-WS index page where the leaves are HTML markup nodes that might be annotated with an attempted RDFa syntax.

%The PT Tree is a three level filesystem /HTML Dom tree of
%Directories per volumes - with an index HTML DOM Tree for the volume and a list of individual PDF files per paper of the volume.

%The leaf node contents are the PDF contents of the files and the HTML markup nodes of the Volume index page.
Both the HTML and PDF encoding of the original scientific content are structured for the purpose of optimizing the display / output on paper or screens; therefore, there is a structure loss compared to what was originally available in the text document processor files that the authors might have been using~\cite{Fahl2023HSP}.  Most scholars are not aware of this loss in the daily use of published content just because the documents are optimized for display and human consumption~\cite{Yu2020}. For the metadata extraction and use in knowledge graphs the difference is sometimes disastrous, e.g., when a simple text from a PDF document can not be extracted any more due to exotic styling and formatting or simply because only a scanned graphic image version of an older document is available that needs optical character recognition to extract the textual content.

The metadata needed for creating the knowledge graph is only available in natural language/text form and follows rules that have been changed multiple times during the history of CEUR-WS. From 2013 to 2026, there have been 33 different versions of the index file template with no proper tracing of what was changed from version to version. The index and volume HTML files were often edited manually, leading to minor, undocumented differences even within these versions. The usage of the different page versions follows a long-tail Zipfian distribution with 5 versions covering 60\% of all volumes.

% https://cr.bitplan.com/index.php/CEUR-WS_Versions
\subsection{Contributions of this Paper}
The core contributions presented in this work are:
\begin{enumerate}
    \item To provide the tools, infrastructure and approaches to to consistently supply high quality FAIR metadata (Semantification) for CEUR-WS as detailed in section ~\ref{sec:Semantification}.

    \item Splitting the metadata for the three main entities Event, Event Series and Proceedings as a major step towards improving the metadata quality and FAIRness.

    \item Providing a bootstrapping~\cite{Bardini2001,Engelbart2007} approach to shift the single point of truth from the current HTML/PDF files to a RDF ready ontology based format as outlined in ~\ref{sec:single_point_of_truth_bootstrapping}.

\end{enumerate}

These contributions are valuable for other publishing use cases that struggle with the same paint points.

\section{From Semantic Publishing Challenge 2015 via Text2KG2023 to Production 2026}\label{sec:SemantificationStory}
The Semantic Publishing Challenge ~\cite{Lange2014,DiIorio2015} started in 2014 with Task~1 calling for \tquote{answering queries related to the quality of workshops by computing metrics from data extracted from their proceedings, also considering information about persons and events}.

The relevance of the original set of 20 queries for Task~1 was subjectively rated by us, resulting in \footnotelink{https://cr.bitplan.com/index.php/List\_of\_Queries}{a list sorting the queries by priority}. The most relevant queries and the 5 queries Q1.5, Q1.12, Q1.13, Q1.16 and Q1.17 that rely on the main index have been taken as the definition of done. We reported on the successful Wikidata export of all CEUR-WS proceeding and event metadata in 2023 at the TEXT2KG Workshop ~\cite{FahlHOD23}.
The \footnotelink{https://cr.bitplan.com/index.php/Semantic\_Publishing\_Challenge\_Queries}{SPARQL queries} are available as named parameterized queries. The snapquery project ~\cite{snapquery} was motivated by this set of queries and was a candidate middleware for the Scholia project ~\cite{Willighagen2026Scholia}. 

https://scholia.toolforge.org/event/Q124669329
https://qlever.scholia.wiki/event/Q124669329

\section{CEUR-WS Semantification in Wikidata}\label{sec:Semantification}

Figure~\ref{fig:ceurmake_workflow} gives an overview of the CEUR-WS Semantification workflow.

\begin{figure}
    \centering
    \includegraphics[width=\linewidth]{images/pyceurmake_workflow.eps}
    \caption{Workflow of the CEUR-WS Semantification}
    \label{fig:ceurmake_workflow}
\end{figure}

For each core entity we analyzed the availability of Wikidata properties by using our wdgrid analysis tool ~\cite{wdgrid} ~\cite{FahlHW0D22}. This tool also generates naive or sophisticated SPARQL queries. We aligned the queries with the stated needs of the Semantic Publishing Challenge. 

The pyCEURmake ~\cite{pyCEURmake} tool was created to parse the main HTML index as well as the individual proceedings HTML pages. The results have been cached in JSON and SQLite storage for quick access. Location entries have been consolidated using the ~\cite3{geograpy3} library. Author, Editor and Paper details have been aligned with the dblp SPARQL endpoint.

snapquery ~\cite{snapquery} was used to make sure the queries are compatible with different endpoints and do not only run on the official blazegraph based Wikidata query endpoint. 

ceur-spt ~\cite{ceurs-spt} is the prototype for showing the benefits of semantification by improved linking and navigation capababilities compared to the current static CEUR-WS HTML page.
 
\section{In-Use: Community Impact - WikiCite, Scholia, CEUR-WS }\label{sec:Community Impact}
WikiCite is an initiative to create a bibliographic database based on Wikidata  ~\cite{WikiCite}. 

Scholia ~\cite{Nielsen2017} is a portal serving wikidata content for the needs of scholars. It is based on a set of over 380 named parameterized SPARQL queries organized around 40 different aspects. The Wikidata Scholarly Articles Subgraph Analysis ~\cite{AKhatun2021} showed that roughly half of the 17 billion triples that form the main graph belong to scholarly articles and their linked entities. The Wikimedia foundation needed to perform a graph split due to a 4 Terrabyte maximum storage limit of the Blazegraph backend in use. The Scholia community decided to use QLever ~\cite{Bast2017} as the new backend and make all queries SPARQL1.1 compliant ~\cite{Willighagen2026Scholia}.
Note that before the graph split there was a hold on adding more scholarly articles to Wikidata which was the reason we did no mass-insert CEUR-WS paper metadata yet.
Writes to Wikidata are performed through the officially approved bot account \textit{User:CEUR-WS}~\cite{CEURWSBot}, for which the community permission was obtained in 2023~\cite{ceur-ws-bot-rquest2023}.

WikiCite, Wikidata, Scholia and the CEUR-WS stakeholder community are overlapping communities. There are over 100.000 distinct authors and over 6.500 distinct editors with digital traces in the dblp CEUR-WS index entries. Since the CEUR-WS wikidata bot 
 
\section{Single Point of Truth Bootstrapping}
\label{sec:single_point_of_truth_bootstrapping}
\section{Lessons Learned}\label{sec:LessonsLearned}
\subsection{Queries as first-class citizens}
\subsection{Property creation as a community process}
Wikidata properties such as \wdp{P11633}{colocated with} undergo a \footnotelink{https://www.wikidata.org/wiki/Wikidata:Property\_proposal}{Wikidata's property proposal process} process before a new propery is created. As described in ~\cite{FahlHOD23} the creation of \footnotelink{https://www.wikidata.org/wiki/Wikidata:Property\_proposal/colocated\_with}{property proposal} took some effort and  lead to the clarification that the property is asymmetric. The property was considered as highly relevant and well defined. 



\subsection{Legal aspects of Publishing Personal Data of Scholars}
\label{sec:legal}
The proposed workflow calls for publishing personal data of authors
and editors early in the life cycle of an event. According to the legal advisors we asked this is legal since a combination of conditions (Art.\ 6 GDPR) is met:
Explicit consent is given by the stakeholders (Art.\ 6(1)(a)), the legitimate interest of the
stakeholders is pursued (Art.\ 6(1)(f)), the data might be public already, the processing is
necessary for scientific research (Art.\ 6(1)(e), Art.\ 89). The strongest condition is that
there might be a legal obligation (Art.\ 6(1)(c)) as would be the case if publishing via major
print publishers that have to provide a copy of the proceedings to the German National Library
according to German law.
Consent, Public Data, and Contractual Necessity are crucial for
scholarly publishing.
\subsection{Iterative introduction beats big bang}
The challenge in introducing the new CEUR-WS workflow is to find a sweet spot between effort and benefit~\cite{Pohl2015-vj}. Given the long-tail distribution of issues of ever more corner and exotic cases that have shown up in the analysis of the existing data and workflow there is potentially a huge list of \tquote{small} problems that would require an inadequately high amount of effort to be solved systematically. The \tquote{sweet spot} is in solving only the most relevant problems and not fixing some of the exotic details at all or doing it manually. Therefore we will avoid a \tquote{big bang} introduction, which would replace  the old approach in one single big (and therefore very risky) step.
Instead we propose to do smaller changes gradually. As a first step, the main index.html shall be split by years (with a link to the old legacy complete list). Then, the per-year list may be generated from the single point of truth data. This will be continued iteratively, going from the current year back to a point where it seems not reasonable any more to spent further effort.
\section{Roadmap: Metadata first Publishing}\label{sec:Metadata-first Publishing}
The coverage of the core entities Event Series, Papers, Scholars and Institutions needs to be improved based on consolidating and completing the RDF data of the downstream indexing services at DBLP and german national library DNB. 
The holy grail of this effort will be to change the CEUR-WS single point of truth from the current HTML/PDF base to a workflow in which the metadata is mined from the proceedings matter. 

\section{Conclusion}\label{sec:Conclusion}
CEUR-WS proceedings metadata is systematically made available in Wikidata via the CEUR-WS bot. The relevant Task~1 Named Parameterized Queries of the Semantic Publishing Challenge of 2014/2015 are now available on different endpoints namely the  SPARQL1.1  compatible QLever endpoint. The Scholia portal provides production access to queries based on and extending this set. The Wikidata, Wikicite, Scholia and CEUR-WS stakeholder communities have appreciated and contributed to this work.

\subsection{Supplemental Material Statement} Source code for the prototypes and semantification tools is available via the repositories of the
\footnotelink{https://github.com/ceurws}{github ceurws organization}. Extended supplemental material --- including the schema diagrams (Figures~\ref{fig:core_entities} and~\ref{fig:swat4hcls_2024_instance}) in their editable Graphviz form, the SPARQL queries derived from the Semantic Publishing Challenge, and background research --- is provided by the CR research wiki~\cite{Fahl_CR_Wiki}.
The LaTeX sources, bibliography, build pipeline and the \texttt{AGENTS.md} governance policy for AI-assisted editing of this paper are available at~\cite{CEURWSSemantificationISWC2026Repo}. The cited software artefacts are released under Apache-2.0: \texttt{ceur-spt}~\cite{ceur-spt}, \texttt{pyCEURmake}~\cite{pyCEURmake}, \texttt{snapquery}~\cite{snapquery}, and \texttt{wdgrid}~\cite{wdgrid}.
\subsection{Acknowledgements}
\label{sec:ack}

We would like to thank Jakob Voß for helping with the k10plus matching and creating the Wikidata property \wdp{colocated with}{P11633} in due time.

Thomas Hoeren and Jonas Kuiter of ITM, Münster kindly gave hints for legal aspects as summarized in section ~\ref{sec:legal}.

This paper is dedicated to the memory of CEUR-WS board member Ralf Klamma $\dag$ January 2023.

This research has been partly funded by a grant of the Deutsche Forschungsgemeinschaft (DFG). \footnote{\href{https://gepris.dfg.de/gepris/projekt/426477583}{ConfIDent project; see https://gepris.dfg.de/gepris/projekt/426477583}}

\subsection{Declaration of use of Generative AI}
The authors used AI-based coding assistants during the engineering and preparation of this work. Use of such assistants is governed by an AGENTS.md policy file maintained in this paper's repository and, equivalently, in the source repositories of the software artefacts ceur-spt\cite{ceur-spt} and pyCEURmake\cite{pyCEURmake}. The policy mandates:
\begin{itemize}
\item a plan-and-confirm protocol: the assistant must state its analysis, plan, and scope, and receive explicit human confirmation before any read, write, or execution;
\item a strict scope for the manuscript itself: AI assistance is limited to spelling/grammar corrections and formatting support (tables, figures, build automation). AI was not used to generate, rewrite or contribute original scientific content, arguments, or prose;
\end{itemize}
All scientific claims, interpretations and conclusions are the sole responsibility of the authors.

% Bibliography
\bibliographystyle{splncs04}
\bibliography{ceurws-semantification}

\end{document}
